%%%%%%%%%%%%%%%%%
% CV tailored for Ingénieur Système Embarqué - Nice
% Positioning: Ingénieur système / intégration / validation | systèmes embarqués critiques, exigences, interfaces HW/SW
%%%%%%%%%%%%%%%%%

\documentclass[8pt,a4paper,ragged2e,withhyper]{altacv}

\geometry{left=1.25cm,right=1.25cm,top=.5cm,bottom=1.cm,columnsep=1.2cm}

\usepackage{paracol}
\usepackage{fontawesome5}
\usepackage{hyperref}

\iftutex 
  \setmainfont{Roboto Slab}
  \setsansfont{Lato}
  \renewcommand{\familydefault}{\sfdefault}
\else
  \usepackage[rm]{roboto}
  \usepackage[defaultsans]{lato}
  \renewcommand{\familydefault}{\sfdefault}
\fi

\definecolor{SlateGrey}{HTML}{2E2E2E}
\definecolor{LightGrey}{HTML}{666666}
\definecolor{DarkPastelRed}{HTML}{8AC0D9}
\definecolor{PastelRed}{HTML}{00B0DA}
\definecolor{GoldenEarth}{HTML}{B4AD0C}

\colorlet{name}{black}
\colorlet{tagline}{PastelRed}
\colorlet{heading}{DarkPastelRed}
\colorlet{headingrule}{GoldenEarth}
\colorlet{subheading}{PastelRed}
\colorlet{accent}{PastelRed}
\colorlet{emphasis}{SlateGrey}
\colorlet{body}{LightGrey}

\definecolor{violet}{HTML}{5A2582}
\definecolor{rouge}{HTML}{F53B3B}
\definecolor{noir}{HTML}{00232C}
\definecolor{bleu}{HTML}{00B0DA}
\definecolor{rose}{HTML}{F831A0}
\definecolor{jaune}{HTML}{FFEC00}
\definecolor{vert}{HTML}{34AD6C}

\renewcommand{\namefont}{\Huge\rmfamily\bfseries}
\renewcommand{\personalinfofont}{\footnotesize}
\renewcommand{\cvsectionfont}{\LARGE\rmfamily\bfseries}
\renewcommand{\cvsubsectionfont}{\large\bfseries}
\renewcommand{\cvItemMarker}{{\small\textbullet}}
\renewcommand{\cvRatingMarker}{\faCircle}

\input{../_shared/pubs-num.tex}
\addbibresource{../_shared/sample.bib}

% auto_tex_manager layout tuning
\emergencystretch=3em
\hfuzz=60pt
\hbadness=9999
\sloppy

\begin{document}

\name{BOILEAU Guillaume}
\tagline{Ingénieur Système / Intégration | Systèmes embarqués critiques, exigences, validation \& interfaces HW/SW}

\photoL{2.8cm}{../_shared/moi}

\personalinfo{%
  \email{guillaume.boileaupro@gmail.com}
  \phone{+33 6 59 94 85 84}
  \location{Alpes-Maritimes, FRANCE}
  \linkedin{guillaume-boileau-891b81265}
  \researchgate{Guillaume-Boileau-2}
  \orcid{0000-0002-3576-6968}
}

\makecvheader
\vspace{-0.3cm}

\AtBeginEnvironment{itemize}{\small}
\columnratio{0.64}

\begin{paracol}{2}

\cvsection{Experience}

\cvevent{\textbf{Software Engineer -- Développement logiciel, intégration \& support opérationnel}}{VEV Platform Services France}{2025 -- 2026}{Valbonne, France}

\textbf{Contexte :} Environnement industriel logiciel pour une plateforme CPMS liée à des infrastructures de recharge de véhicules électriques, avec services backend, interfaces applicatives, données opérationnelles et équipements terrain.

\textbf{Objectif :} Contribuer au développement, à l’intégration et à la fiabilisation de services logiciels connectés, en analysant les flux de données, les incidents et les comportements applicatifs en environnement opérationnel.

\begin{itemize}
  \item \textbf{Interfaces logiciel / terrain :} Travail sur les interfaces entre services applicatifs, données opérationnelles, infrastructures connectées et contraintes terrain.
  \item \textbf{Intégration logicielle :} Contribution à l’intégration de composants applicatifs dans une plateforme opérationnelle connectée à des équipements.
  \item \textbf{Analyse d’incidents :} Investigation d’interruptions de flux, incohérences de données, comportements inattendus et impacts sur le service.
  \item \textbf{Validation opérationnelle :} Vérification de la cohérence entre données applicatives, comportement des équipements et attendus fonctionnels.
  \item \textbf{Développement logiciel :} Contribution à des composants TypeScript, Angular et Node.js.
  \item \textbf{Documentation technique :} Formalisation de constats, analyses d’anomalies et éléments utiles à la résolution par les équipes techniques.
\end{itemize}

\divider

\cvevent{\textbf{Ingénieur R\&D senior -- Ingénierie système, validation \& performance spatiale}}{CNES}{2023 -- 2025}{Nice, France}

\textbf{Contexte :} Activités d’ingénierie et de R\&D pour la mission spatiale LISA ESA/NASA, dans un environnement international impliquant simulation, développement Python, traitement de données, intégration de modèles et validation technique.

\textbf{Objectif :} Structurer, intégrer, comparer et valider des contributions techniques complexes afin de produire des résultats exploitables, traçables et présentables en revue collaborative.

\begin{itemize}
  \item \textbf{Ingénierie système :} Analyse de chaînes techniques complexes, identification d’hypothèses, interfaces, contraintes et dépendances entre contributions.
  \item \textbf{Exigences \& traçabilité :} Structuration de configurations, versions, données d’entrée, sorties de simulation et résultats d’analyse pour assurer la traçabilité.
  \item \textbf{Validation système :} Définition de contrôles de cohérence, règles de comparaison, méthodes de vérification et critères d’acceptabilité des résultats.
  \item \textbf{Analyse de performance :} Évaluation de résultats de simulation dans des contextes de signaux faibles, bruit, incertitudes et contraintes de cohérence.
  \item \textbf{Analyse d’anomalies :} Identification d’écarts entre modèles, analyse d’impact, formalisation des points bloquants et proposition d’actions correctives.
  \item \textbf{Coordination technique :} Interface avec des contributeurs scientifiques, logiciels et institutionnels dans un environnement spatial international.
  \item \textbf{Revues techniques :} Consolidation de livrables, synthèses techniques, traçabilité des hypothèses et préparation d’éléments de décision.
  \item \textbf{Plateforme Python :} Développement de CATpyLISA pour l’intégration, la comparaison, la validation et la revue technique de modèles.
  \textcolor{DarkPastelRed}{\href{https://gitlab.oca.eu/gboileau/lisa_l3_tools}{\bf GitLab: CATpyLISA}}
\end{itemize}

\divider

\cvevent{\textbf{Data Scientist -- Données capteurs, bruit \& analyse de performance}}{Universiteit Antwerpen, Belgium}{2021 -- 2023}{Antwerp, Belgium}

\textbf{Contexte :} Études de performance pour instruments de précision et détecteurs de nouvelle génération, combinant mesures environnementales, bruit, modélisation statistique, configurations instrumentales et validation de résultats.

\textbf{Objectif :} Développer des workflows robustes et reproductibles pour identifier, caractériser et réduire des sources affectant la qualité de mesure et la sensibilité de systèmes complexes.

\begin{itemize}
  \item \textbf{Données capteurs :} Analyse de mesures issues de capteurs environnementaux, corrélations spatiales et effets de propagation.
  \item \textbf{Analyse de performance :} Évaluation de la qualité de mesure, des limites instrumentales et de l’impact de différentes sources de bruit.
  \item \textbf{Traitement du signal :} Identification et caractérisation de contributions de bruit dans des mesures complexes.
  \item \textbf{Configurations système :} Études de différentes configurations instrumentales et de différents niveaux de bruit pour la détection de signaux très ténus.
  \item \textbf{Validation de résultats :} Analyse de sensibilité, vérification de cohérence, figures de mérite et évaluation des limites de performance.
  \item \textbf{Reporting technique :} Production de résultats traçables, de synthèses techniques et de supports pour parties prenantes internationales.
\end{itemize}

\divider

\cvevent{\textbf{Ingénieur R\&D junior -- Signal, simulation \& diagnostics techniques}}{ARTEMIS, Observatoire de la Côte d’Azur}{2018 -- 2021}{Nice, France}

\textbf{Contexte :} Activités de R\&D liées à l’analyse de signaux faibles, à la simulation numérique, aux diagnostics de bruit et à la préparation de missions spatiales et d’instruments de détection.

\textbf{Objectif :} Développer des méthodes d’analyse, automatiser des simulations et produire des résultats exploitables dans un environnement scientifique et technique exigeant.

\begin{itemize}
  \item \textbf{Simulation numérique :} Développement d’approches de simulation pour environnements de signaux complexes et jeux de données de grande taille.
  \item \textbf{Analyse de données :} Mise en place de méthodes pour analyser et séparer des signaux faibles et superposés.
  \item \textbf{Traitement du signal :} Études de caractérisation, de décomposition et de comparaison de signaux dans des contextes bruités.
  \item \textbf{Validation technique :} Vérification de cohérence, comparaison de résultats et études de sensibilité.
  \item \textbf{Diagnostics :} Analyse de sources de bruit, corrélations entre capteurs et risques d’interprétation des résultats.
  \item \textbf{Documentation :} Formalisation des méthodes, hypothèses, résultats et conclusions pour revue collaborative.
\end{itemize}

\switchcolumn

\cvsection{Profile}

\textbf{Ingénieur docteur orienté ingénierie système, intégration, validation et analyse de performance, avec une expérience solide en systèmes spatiaux, environnements complexes, données capteurs, traitement du signal, logiciel et documentation technique.}

Positionnement pour ce rôle : ingénieur système capable d’intervenir sur la cohérence d’architectures multi-domaines, la déclinaison d’exigences, l’intégration de sous-ensembles, la validation et la qualification, avec une montée ciblée sur calculateurs embarqués durcis, interfaces HW/SW/mécanique et contraintes environnementales.

\begin{itemize}
  \item Expérience en validation système, analyse d’écarts, traçabilité, documentation technique et préparation de revues.
  \item Culture forte des systèmes complexes, du spatial, de la simulation, des données capteurs et de l’analyse de performance.
  \item Capacité à dialoguer avec des équipes logiciel, système, instrumentation, qualité et métiers techniques.
  \item Profil rigoureux, autonome, orienté exigences, conformité, diagnostic technique et résolution de problèmes.
\end{itemize}

\cvsection{Languages}

\cvskill{English}{4}
\divider
\cvskill{Spanish}{2}
\divider

\medskip

\cvsection{Education}

\cvevent{PhD in Physics}{Université Côte d’Azur}{2018 -- 2021}{Nice, France}

\divider

\cvevent{Selective Graduate Program in Fundamental Physics (Magistère)}{Université Paris-Saclay}{2016 -- 2018}{Orsay, France}

\divider

\cvevent{MSc in Astronomy and Astrophysics}{Observatoire de Paris / Université Paris-Saclay}{2017 -- 2018}{Paris, France}

\divider

\cvevent{BSc in Fundamental Physics}{Université Paris-Saclay}{2015 -- 2016}{Orsay, France}

\cvsection{Professional Training}

\cvevent{Machine Learning in Python with scikit-learn}{FUN MOOC -- Inria}{2026}{France Université Numérique}

\small{
Predictive modelling, supervised learning, model evaluation, cross-validation, metrics, pipelines and practical machine learning workflows with Python and scikit-learn.
}

\divider

\cvevent{Supervised Machine Learning: Regression \& Classification}{DeepLearning.AI -- Andrew Ng}{2020}{Coursera}

\divider

\cvevent{Recherche reproductible : principes \newline méthodologiques pour une science transparente}{FUN MOOC -- Inria}{2025}{France Université Numérique}


\divider

\cvevent{Continuous Professional Training -- Software, Data, AI and Systems}{OpenClassrooms}{2024 -- 2026}{}

\small{
Python, SQL, Machine Learning, AI, Linux, JavaScript, TypeScript, Angular, Node.js, Django, REST APIs, C/C++, web development, algorithms and software engineering fundamentals.
}

\end{paracol}

\cvsection{Projects}

\cvevent{CATpyLISA -- Plateforme Python d’intégration, comparaison et validation}{Mission spatiale LISA ESA/NASA}{2023 -- 2025}{}

Plateforme Python dédiée à l’intégration de modèles, à la structuration de données, à la comparaison de résultats, à la validation technique et à la revue collaborative de workflows scientifiques.

\textbf{Rôle :} développement Python, structuration de données, documentation technique, validation, analyse d’écarts, synthèse de résultats, coordination avec plusieurs contributeurs.

\textbf{Compétences mobilisées :} Python, intégration de modèles, validation système, traitement du signal, analyse de performance, simulation, documentation, GitLab, Confluence, environnement spatial international.

\divider

\cvevent{Workflows de simulation, performance et analyse de bruit}{Instrumentation scientifique et préparation de missions spatiales}{2018 -- 2025}{}

Développement de workflows numériques et statistiques pour analyser des signaux faibles, caractériser des contributions de bruit, comparer des résultats et évaluer la robustesse de modèles dans des environnements complexes.

\textbf{Rôle :} analyse de performance, simulation, traitement du signal, caractérisation de bruit, études de sensibilité, diagnostics techniques, reporting et support à la décision.

\textbf{Compétences mobilisées :} traitement du signal, simulation numérique, analyse spectrale, Python, C, validation, analyse de sensibilité, documentation technique.

\divider

\cvevent{VEV -- Développement logiciel, intégration \& support opérationnel}{Industrial software / Data flows / Connected infrastructure}{2025 -- 2026}{}

Développement et intégration de services logiciels pour une plateforme CPMS connectée à des infrastructures de recharge, avec analyse de flux FTP, investigation d’incidents, vérification de comportements applicatifs, support opérationnel et documentation technique.

\textbf{Rôle :} développement TypeScript, Angular et Node.js, intégration logicielle, analyse de flux FTP, investigation d’incidents, vérification de comportements applicatifs et formalisation de constats techniques.

\textbf{Compétences mobilisées :} TypeScript, Angular, Node.js, FTP, Linux, Git, analyse d’incidents, intégration logicielle, support opérationnel, documentation technique.

\cvsection{Compétences clés}

\vspace{-.3cm}

\cvsubsection{Ingénierie système, exigences \& conformité}

{\small
\cvtag{Ingénierie système}
\cvtag{Analyse d’exigences}
\cvtag{Déclinaison d’exigences}
\cvtag{Traçabilité}
\cvtag{Documentation système}
\cvtag{Conformité qualité}
\cvtag{EN9100 -- awareness}
\cvtag{AQAP -- awareness}
\cvtag{Revues techniques}
\cvtag{Support décisionnel}
\cvtag{Systèmes critiques}
\cvtag{Défense -- habilitation compatible}
}

\divider\smallskip
\vspace{-.3cm}

\cvsubsection{Intégration, validation \& qualification}

{\small
\cvtag{Intégration système}
\cvtag{Intégration de sous-ensembles}
\cvtag{Validation système}
\cvtag{IVVQ}
\cvtag{Qualification -- ramp-up}
\cvtag{CEM -- ramp-up}
\cvtag{Thermique -- awareness}
\cvtag{Vibrations -- awareness}
\cvtag{Chocs -- awareness}
\cvtag{Étanchéité -- awareness}
\cvtag{Procédures de test}
\cvtag{Analyse de résultats}
\cvtag{Traitement d’anomalies}
}

\divider\smallskip
\vspace{-.3cm}

\cvsubsection{Systèmes embarqués, HW/SW \& interfaces}

{\small
\cvtag{Systèmes embarqués -- ramp-up}
\cvtag{Calculateurs embarqués -- ramp-up}
\cvtag{Calculateurs durcis -- awareness}
\cvtag{Interfaces hardware/software}
\cvtag{Architecture système}
\cvtag{Cartes électroniques -- awareness}
\cvtag{CPU / GPU -- awareness}
\cvtag{FPGA -- awareness}
\cvtag{Interfaces mécaniques -- awareness}
\cvtag{Bus de communication -- ramp-up}
\cvtag{Données capteurs}
\cvtag{Instrumentation scientifique}
}

\divider\smallskip
\vspace{-.3cm}

\cvsubsection{Analyse, diagnostic \& performance}

{\small
\cvtag{Analyse d’anomalies}
\cvtag{Root cause analysis}
\cvtag{Diagnostic technique}
\cvtag{Analyse de performance}
\cvtag{Contrôles de cohérence}
\cvtag{Validation de résultats}
\cvtag{Analyse de sensibilité}
\cvtag{Simulation numérique}
\cvtag{Modélisation}
\cvtag{Traitement du signal}
\cvtag{Analyse de bruit}
\cvtag{Systèmes complexes}
}

\divider\smallskip
\vspace{-.3cm}

\cvsubsection{Logiciel, data \& outils}

{\small
\cvtag{Python}
\cvtag{C}
\cvtag{Pandas}
\cvtag{NumPy}
\cvtag{SciPy}
\cvtag{Matplotlib}
\cvtag{SQL}
\cvtag{TypeScript}
\cvtag{Angular}
\cvtag{Node.js}
\cvtag{Bash}
\cvtag{Linux}
\cvtag{Git}
\cvtag{GitLab}
\cvtag{Docker}
\cvtag{CI/CD}
\cvtag{Confluence}
\cvtag{Jira}
\cvtag{Zenhub}
\cvtag{OpenSearch}
\cvtag{Sentry}
}

\end{document}
